\section{Introduction}

In autonomous driving, coverage analysis is a crucial step to ensure the safety and reliability of the system.
At its core, a coverage argument asks whether the scenes observed so far have exhausted the space of traffic scenes that can occur in the real world.
In most situations, such arguments are collected either per coverage factor, or maybe up to 2 or 3 factor interactions.
See for example \cite{foretellix2019} for a production-grade implementation of state-of-the-art coverage analysis.

In contrast to existing approaches, this paper proposes a graph-based framework for coverage analysis.
The framework takes the \emph{complete traffic scene} as the unit of coverage rather than an individual factor: each scene is represented as a single graph whose nodes carry the parameters describing the individual actors and whose edges carry the parameters describing the relations between them.
The key advantage of this representation is that arbitrarily complex scenes---with varying numbers of actors and the higher-order interactions between them---are captured naturally within one object, instead of being decomposed into isolated coverage factors or low-order factor combinations.
Coverage is therefore assessed against whole scenes: we ask how completely the observed scene graphs cover those that can occur in reality.
While graph-based representations of traffic scenes already exist, they have not been specifically designed for this coverage analysis.
This work bridges that gap by utilizing graph-based traffic scene representations explicitly for coverage analysis.

This paper is concerned with completeness of the observed scenes, not with the safety of the system. A reasonable strategy is to focus coverage analysis on the more safety-critical scenes, but here the focus is solely on coverage, and we leave safety considerations to future work.


